% GENERATED FILE. Edit canonical lessons, then run npm run generate:books.
% canonical-source-hash: fnv1a64:1b041a062a37978e
% canonical-lessons: ML-C137-vehicle, ML-C137-car, ML-C137-train, ML-C137-plane, ML-C137-auto

\chapter{Getting There}
\label{ch:ml-getting-there}
\hlchaptermodality{137}

\begin{chapteropening}
\textbf{By the end of this chapter:} \emph{I can name a vehicle, a car, a train, a plane and an autorickshaw.}

Complete the last lesson of chapter 137: I can name a vehicle, a car, a train, a plane and an autorickshaw.
\end{chapteropening}

\section[\texorpdfstring{\ml{വണ്ടി} (vaṇṭi)}{വണ്ടി (vaṇṭi)}]{\ml{വണ്ടി} (vaṇṭi) — a vehicle}

\label{lesson:ML-C137-vehicle}



\begin{quote}
Before the new one: say the Malayalam for a well, then the Malayalam for a harbour.
\end{quote}

\subsection*{You'll want to know: \ml{വണ്ടി}}
\textbf{\ml{വണ്ടി}} (\emph{vaṇṭi}) — \textquotedblleft{}a vehicle\textquotedblright{}, any cart, car or bus.

\begin{cousinweb}
\textbf{\ml{വണ്ടി}} (\emph{vaṇṭi}) is listed in DEDR appendix 50, among Dravidian words taken from Indo-Aryan, with Tamil \emph{vaṇṭi} and Kannada \emph{baṇḍi}, \textquotedblleft{}cart\textquotedblright{}. DEDR traces the word ultimately to Sanskrit \emph{bhāṇḍa}, \textquotedblleft{}goods, wares\textquotedblright{} (CDIAL 9440): the cart was named for what it carried.
\end{cousinweb}

\begin{grammarlens}[title={What you've built}]
The general word.
\end{grammarlens}

\subsection*{Your turn}
\begin{itemize}
\raggedright
  \item \emph{Say it:} \emph{vaṇṭi}
  \item \emph{Say it:} \emph{vaṇṭi}, once more
  \item \emph{Say it:} say \emph{vaṇṭi}
  \item \emph{Recall:} say the Malayalam for a well, then the Malayalam for a harbour, then say \emph{vaṇṭi} again
\end{itemize}

\begin{tcolorbox}[breakable,colback=teal!4,colframe=teal!35!black,title={Before you move on}]
What does \ml{വണ്ടി} mean? (\textquotedblleft{}A vehicle\textquotedblright{}.) Say it once more.
\end{tcolorbox}
\section[\texorpdfstring{\ml{കാർ} (kār)}{കാർ (kār)}]{\ml{കാർ} (kār) — a car}

\label{lesson:ML-C137-car}



\begin{quote}
Before the new one: say the Malayalam for a harbour, then the Malayalam for a vehicle.
\end{quote}

\subsection*{You'll want to know: \ml{കാർ}}
\textbf{\ml{കാർ}} (\emph{kār}) — \textquotedblleft{}a car\textquotedblright{}, from English.

\begin{cousinweb}
\textbf{\ml{കാർ}} (\emph{kār}) $\leftarrow$ English \emph{car}, which goes back to Latin \emph{carrus}, \textquotedblleft{}a wagon\textquotedblright{}.
\end{cousinweb}

\begin{grammarlens}[title={What you've built}]
A private vehicle.
\end{grammarlens}

\subsection*{Your turn}
\begin{itemize}
\raggedright
  \item \emph{Say it:} \emph{kār}
  \item \emph{Say it:} \emph{kār}, once more
  \item \emph{Say it:} say \emph{kār}
  \item \emph{Recall:} say the Malayalam for a harbour, then the Malayalam for a vehicle, then say \emph{kār} again
\end{itemize}

\begin{tcolorbox}[breakable,colback=teal!4,colframe=teal!35!black,title={Before you move on}]
What does \ml{കാർ} mean? (\textquotedblleft{}A car\textquotedblright{}.) Say it once more.
\end{tcolorbox}
\section[\texorpdfstring{\ml{തീവണ്ടി} (tīvaṇṭi)}{തീവണ്ടി (tīvaṇṭi)}]{\ml{തീവണ്ടി} (tīvaṇṭi) — a train}

\label{lesson:ML-C137-train}



\begin{quote}
Before the new one: say the Malayalam for a vehicle, then the Malayalam for a car.
\end{quote}

\subsection*{You'll want to know: \ml{തീവണ്ടി}}
\textbf{\ml{തീവണ്ടി}} (\emph{tīvaṇṭi}) — \textquotedblleft{}a train\textquotedblright{}, literally \textquotedblleft{}fire-vehicle\textquotedblright{}: \textbf{\ml{തീ}} (\emph{tī}), \textquotedblleft{}fire\textquotedblright{}, with \textbf{\ml{വണ്ടി}}.

\begin{cousinweb}
\textbf{\ml{തീവണ്ടി}} (\emph{tīvaṇṭi}) is \textquotedblleft{}the fire-cart\textquotedblright{}: native \emph{tī}, \textquotedblleft{}fire\textquotedblright{} (DEDR 3266, with Tamil \emph{tī}), joined to \emph{vaṇṭi}, \textquotedblleft{}cart\textquotedblright{} (DEDR appendix 50, with Tamil \emph{vaṇṭi}).
\end{cousinweb}

\begin{grammarlens}[title={What you've built}]
A vehicle word you can take apart.
\end{grammarlens}

\subsection*{Your turn}
\begin{itemize}
\raggedright
  \item \emph{Say it:} \emph{tīvaṇṭi}
  \item \emph{Say it:} \emph{tīvaṇṭi}, once more
  \item \emph{Say it:} say \emph{tīvaṇṭi}
  \item \emph{Recall:} say the Malayalam for a vehicle, then the Malayalam for a car, then say \emph{tīvaṇṭi} again
\end{itemize}

\begin{tcolorbox}[breakable,colback=teal!4,colframe=teal!35!black,title={Before you move on}]
What does \ml{തീവണ്ടി} mean? (\textquotedblleft{}A train\textquotedblright{}.) Say it once more.
\end{tcolorbox}
\section[\texorpdfstring{\ml{വിമാനം} (vimānaṁ)}{വിമാനം (vimānaṁ)}]{\ml{വിമാനം} (vimānaṁ) — a plane}

\label{lesson:ML-C137-plane}



\begin{quote}
Before the new one: say the Malayalam for a car, then the Malayalam for a train.
\end{quote}

\subsection*{You'll want to know: \ml{വിമാനം}}
\textbf{\ml{വിമാനം}} (\emph{vimānaṁ}) — \textquotedblleft{}an aeroplane\textquotedblright{}.

\begin{cousinweb}
\textbf{\ml{വിമാനം}} (\emph{vimānaṁ}) $\leftarrow$ Sanskrit \emph{vimāna}, \textquotedblleft{}measure; a car of the gods\textquotedblright{}, taken over whole as a tatsama, with Malayalam's noun ending \emph{-ṁ} (CDIAL 11823). The flying chariot of the old stories gave its name to the aeroplane.
\end{cousinweb}

\begin{grammarlens}[title={What you've built}]
The way to Kerala from far away.
\end{grammarlens}

\subsection*{Your turn}
\begin{itemize}
\raggedright
  \item \emph{Say it:} \emph{vimānaṁ}
  \item \emph{Say it:} \emph{vimānaṁ}, once more
  \item \emph{Say it:} say \emph{vimānaṁ}
  \item \emph{Recall:} say the Malayalam for a car, then the Malayalam for a train, then say \emph{vimānaṁ} again
\end{itemize}

\begin{tcolorbox}[breakable,colback=teal!4,colframe=teal!35!black,title={Before you move on}]
What does \ml{വിമാനം} mean? (\textquotedblleft{}A plane\textquotedblright{}.) Say it once more.
\end{tcolorbox}
\section[\texorpdfstring{\ml{ഓട്ടോ} (ōṭṭō)}{ഓട്ടോ (ōṭṭō)}]{\ml{ഓട്ടോ} (ōṭṭō) — an autorickshaw}

\label{lesson:ML-C137-auto}



\begin{quote}
Before the new one: say the Malayalam for a train, then the Malayalam for an aeroplane.
\end{quote}

\subsection*{You'll want to know: \ml{ഓട്ടോ}}
\textbf{\ml{ഓട്ടോ}} (\emph{ōṭṭō}) — \textquotedblleft{}an autorickshaw\textquotedblright{}, the three-wheeled taxi.

\begin{cousinweb}
\textbf{\ml{ഓട്ടോ}} (\emph{ōṭṭō}) $\leftarrow$ English \emph{auto}, short for \emph{autorickshaw}; \emph{auto-} is Greek for \textquotedblleft{}self\textquotedblright{}, and \emph{rickshaw} is from Japanese \emph{jinrikisha}, \textquotedblleft{}man-power vehicle\textquotedblright{}.
\end{cousinweb}

\begin{grammarlens}[title={What you've built}]
The short ride across town.
\end{grammarlens}

\subsection*{Your turn}
\begin{itemize}
\raggedright
  \item \emph{Say it:} \emph{ōṭṭō}
  \item \emph{Say it:} \emph{ōṭṭō}, once more
  \item \emph{Say it:} say \emph{ōṭṭō}
  \item \emph{Recall:} say the Malayalam for a train, then the Malayalam for an aeroplane, then say \emph{ōṭṭō} again
\end{itemize}

\begin{tcolorbox}[breakable,colback=teal!4,colframe=teal!35!black,title={Before you move on}]
What does \ml{ഓട്ടോ} mean? (\textquotedblleft{}An autorickshaw\textquotedblright{}.) Say it once more.
\end{tcolorbox}
